\section{Radar Emitters and Operating Modes}\label{sec:bg:emitters}

In passive \gls{esm} and \gls{elint}, a \emph{radar emitter} denotes a physical transmitting source whose emissions can be tied to a train of pulses in time and frequency \parencite{wiley2006elint}, though this coherence has become increasingly complex to discern in modern, modulated radar systems.

An \emph{operating mode}, by contrast, is a parameterized configuration of the transmitted waveform and scheduling policy: nominal carrier frequency or hopping law, pulse width, repetition structure, beam scan program, and any agility rules that govern how those quantities change from pulse to pulse \parencite{skolnik2008radar,wiley2006elint}.
Modes are therefore properties of the emission process at a given time, whereas emitters are stable identities of hardware and platform; the same hardware cycles through modes as the mission evolves.

Illustrative mode families include wide-area search, acquisition and track, fire-control illumination, weather sensing, and terminal guidance support \parencite{skolnik2008radar}.
Each family implies characteristic constraints on \gls{pri}/\gls{prf}, dwell time, frequency agility, and antenna motion: search modes favor broad coverage and may use staggered or agile repetition to balance range and velocity ambiguity, while dedicated tracking or illuminator modes often concentrate energy in fewer directions with higher update rates on a target of interest.
Navigation radars on ships and aircraft emphasize short-range obstacle and surface mapping; their emissions tend to repeat with comparatively regular timing tied to a fixed sector or rotating beam pattern, whereas long-range surveillance and fire-control systems interleave search and track waveforms and may switch \gls{rf} bands or burst structures when countering interference \parencite{skolnik2008radar,wiley2006elint}.
The resulting diversity matters for downstream analysis because operators must interpret not only ``which radar,'' but also ``what it is doing now,'' and because the same measured \gls{pdw} parameters can arise from different combinations of hardware and scheduling (\cref{sec:bg:pdw}).

The relationship between emitters and modes is many-to-many.
A single emitter implements many modes over a mission timeline, so a time-ordered pulse stream is a mixture of mode segments whose boundaries need not align with changes of physical source.
Conversely, different emitters can realize the same functional mode with different parameter sets---for example, two fire-control radars may both run a track waveform but with distinct nominal \glspl{prf} and frequency plans.
That structure motivates joint reasoning about physical identity and instantaneous configuration: clustering only over pulses loses scheduling context, whereas treating modes without emitter linkage obscures which trains belong together after deinterleaving errors or overlap \parencite{wiley2006elint}.
The remainder of this chapter situates passive reception in the EW/ESM processing chain (\cref{sec:bg:elint_pipeline}) and then supplies the learning background needed to formalize such joint structure in the method.
